% !TeX root = ../colloquium-2.tex

% Определение первообразной. 
% Формулировка основного свойства первообразных.

\phantomsection\label{sec_Definition20}
\begin{definition}[
    Определение первообразной.
    Формулировка основного свойства первообразных.
]
    \subsection*{Определение первообразной.}
    \begin{note}[purple]
        Пусть функции $F$ и $f$ определены на интервале $(a, \; b)$. Если для любого числа
        $x \in (a, \; b)$ $F'(x) = f(x)$, то функция $F$ называется (\textit{точной})
        первообразной функции $f$ на интервале $(a, \; b)$. 
    \end{note}

    \subsection*{Формулировка основного свойства первообразных.}
    \begin{note}[purple]
        Пусть функции $F_1$ и $F_2$ являются первообразными функции $f$ на интервале $(a, \; b)$.
        Тогда существует такое число $C \in \mathbb{R}$, что $F_1(x) - F_2(x) = C \; 
        \forall x \in (a, \; b)$ 
    \end{note}
\end{definition}